\subsection{REST API}
\label{subsec:github_api_rest}
\REST \API umožňuje přístup k~datům a~funkcím GitHubu prostřednictvím \HTTP požadavků. Každý požadavek na \REST \API obsahuje \HTTP metodu a~cestu. V~závislosti na konkrétním endpointu musí být specifikovány další parametry, například hlavičky, informace o~autentizaci, parametry dotazu nebo tělo požadavku~\cite{github_rest_api}.
Pro interakci s~různými zdroji jsou využívány základní \HTTP metody~\cite{github_rest_api}:
\begin{itemize}
  \item \code{GET} pro získání dat (např. seznam issues nebo detail \PR),
  \item \code{POST} pro vytvoření nových zdrojů (např. vytvoření issue),
  \item \code{PUT} pro aktualizaci existujících zdrojů,
  \item \code{DELETE} pro odstranění zdrojů,
  \item \code{PATCH} pro částečnou aktualizaci zdrojů (např. změna stavu issue).
\end{itemize}
Každý endpoint je definován cestou. Cesta může obsahovat statické i~dynamické části. Například endpoint pro získání issues repozitáře je \code{GET /repos/\{owner\}/\{repo\}/issues}, kde \code{owner} a~\code{repo} představují dynamické části, které musí být nahrazeny konkrétními hodnotami~\cite{github_rest_api}. Ukázka takového požadavku je uvedena ve výpisu~\ref{lst:github_rest_example}.

\begin{listing}[H]
\begin{minted}{shell}
curl --request GET \
--url "https://api.github.com/events?per_page=2&page=1" \
--header "Accept: application/vnd.github+json" \
--header "X-GitHub-Api-Version: 2022-11-28" \
  https://api.github.com/events
\end{minted}
  \caption{Ukázka \REST \API požadavku na GitHub pro získání událostí celého GitHubu.}
  \label{lst:github_rest_example}
\end{listing}

\begin{listing}[H]
\begin{minted}[linenos=false]{text}
HTTP/2 200 OK
Content-Type: application/json; charset=utf-8
\end{minted}
\begin{minted}{json}
[
  {
    "id": "11056651313",
    "type": "PushEvent",
    "actor": {
      "id": 41898282,
      "login": "github-actions[bot]",
      "display_login": "github-actions",
      "url": "https://api.github.com/users/github-actions[bot]"
    },
    "repo": {
      "id": 1184208499,
      "name": "IsBenben/xxxh-fans",
      "url": "https://api.github.com/repos/IsBenben/xxxh-fans"
    },
    "payload": {
      "repository_id": 1184208499,
      "push_id": 33348345827
    },
    "public": true,
    "created_at": "2026-04-26T09:11:54Z"
  }
]
\end{minted}
  \caption{Ukázka odpovědi \REST \API na požadavek ze výpisu~\ref{lst:github_rest_example} (zkráceno).}
  \label{lst:github_rest_response}
\end{listing}

Data jsou nejčastěji zprostředkována ve formátu \JSON, který zajišťuje dobrou čitelnost jak pro člověka, tak pro strojové zpracování. V~závislosti na endpointu lze využít také jiné typy formátování obsahu, například \code{diff} nebo \code{patch}~\cite{github_rest_api}.

GitHub \REST \API je verzováno a~využívá kalendářní verzování. Hlavička GitHub \API verze \code{X-GitHub-Api-Version} není povinná, ale je silně doporučována. Pokud tato hlavička není uvedena, je použita výchozí verze definovaná v~dokumentaci, například \code{2022-11-28}. Verze \API se do požadavku zapisuje ve formátu \code{X-GitHub-Api-Version: 2022-11-28}. Tím je zajištěno, že novější verze nezmění chování integrace navázané na starší verzi.
GitHub také poskytuje seznam zpětně nekompatibilních změn mezi verzemi, návody a dokumentaci pro migraci na novější verze \API pro jednodušší přechod~\cite{github_rest_api_versions}.

I~když jsou některé endpointy veřejně dostupné bez autentizace, většina z~nich vyžaduje autentizaci pomocí tokenu v~hlavičce požadavku. Token lze získat více způsoby, například pomocí \PAT, prostřednictvím GitHub App nebo jako vestavěný \code{GITHUB\_TOKEN} v~rámci GitHub Actions~\cite{github_actions_docs,github_rest_api}.
